\subsection{The primary Hall factor in the level projector}

Let $k=\mathbb F_q$. In this subsection $P$ denotes a monic irreducible
polynomial over $k$, with $P(0)\ne0$, of degree $d\ge1$.
Fix an injective character $T:\overline{k}^{\times}\to\mathbb C^{\times}$
and a root $\alpha$ of $P$ in $k_d=\mathbb F_{q^d}$. Write
$T(\alpha)=\exp(2\pi i a/(q^d-1))$ and put $\eta=T^a$.
On $k_d^\times$ the result depends only on $a\bmod (q^d-1)$.
The assignment $\alpha\mapsto\eta|_{k_d^\times}$ is an isomorphism
from a cyclic group to its dual and commutes with the $q$-power map.
Thus $\eta|_{k_d^\times}$ is regular, meaning that its Frobenius orbit
has size $d$.

Here is the Hall formula defining the primary factor. For
$g\in\operatorname{GL}_d(k)$, let $V=k^d$ be the $k[X]$-module with
$X$ acting as $g$. Write $\mu(0,W)$ for the M\"obius function of its
submodule lattice. If $\beta_\eta(h)$ sums $\eta$ on the eigenvalues
of $h$, with algebraic multiplicities, then
\[
 \tau_P(g)=(-1)^d\sum_{W\subseteq V\,\text{a submodule}}
       \beta_\eta(g|_W)\mu(0,W).
\]
Indeed the Hall product with the constant function one sums over invariant
subspaces, and its Hall inverse evaluated on a module gives
the M\"obius function by the defining recursion. These are the
\lineSrc{Lib/LocalGL/Green/HallAlgebra}{549}{eigenvalue sum and Hall product}
and the \lineSrc{Lib/LocalGL/Green/TauFunctions}{133}{rank and sign convention}
used in the projector.

\begin{lemma}\label{lem:fitted-primary-character}
The function $\tau_P$ on $\operatorname{GL}_d(k)$ is the cuspidal
character attached to the orbit of $\eta|_{k_d^\times}$. Its value
at the identity is $\prod_{j=1}^{d-1}(q^j-1)$.
\end{lemma}
\begin{proof}
Only semisimple $W$ contribute. Indeed, for a nonsemisimple module $W$
its socle, the sum of all its simple submodules, is nonzero and proper.
Induction on length gives vanishing for each proper nonsemisimple
submodule. The sum of the M\"obius values on the other proper submodules
equals the sum over the entire submodule lattice of the socle, hence is
zero. The M\"obius recursion gives $\mu(0,W)=0$.

For each distinct monic irreducible factor $H$ of the characteristic
polynomial, let $s_H=\deg H$, $Q_H=q^{s_H}$, and $m_H$ be the multiplicity
of $k[X]/(H)$ in the socle of $V$. Each nonzero semisimple submodule is
specified by subspaces of $\mathbb F_{Q_H}^{m_H}$. Write $n_H$ for their
dimensions and $B_H=\sum_{H(z)=0}\eta(z)$, summing over distinct roots.
Then $\beta_\eta(g|_W)=\sum_H n_H B_H$. For a vector space of
dimension $n$ over $\mathbb F_Q$, its M\"obius value is
$(-1)^n Q^{\binom n2}$. This follows from the M\"obius recursion and
the finite $q$-binomial identity
\[
 \sum_{n=0}^{m}{m\brack n}_Q(-1)^n Q^{\binom n2}u^n
       =\prod_{j=0}^{m-1}(1-Q^j u),
\]
where the Gaussian coefficient counts the $n$-dimensional subspaces.
M\"obius values multiply on products of lattices. It follows that the
unsigned Hall sum before multiplication by $(-1)^d$ is
\[
 \left.
 \left(\sum_H B_H u_H\frac{\partial}{\partial u_H}\right)
 \prod_H\prod_{j=0}^{m_H-1}(1-Q_H^j u_H)
 \right|_{u_H=1\ \mathrm{for\ all}\ H}.
\]
With two or more distinct factors every term vanishes. With one factor $H$
the value of $\tau_P(g)$ is
\[
       (-1)^{d-1}
       \left(\sum_{H(z)=0}\eta(z)\right)
              \prod_{j=1}^{m_H-1}(1-q^{s_H j}).
\]
In this case $s_H\mid d$, so the roots occurring in the sum lie in $k_d$.
Moreover $m_H=\dim_{k_d}\ker(g-z)$ after extending scalars, as is seen
on each cyclic primary summand. These values on all conjugacy classes agree
with Green's cuspidal character formula for the regular character
$\eta|_{k_d^\times}$, in the form stated by D.~Prasad,
\emph{The space of degenerate Whittaker models for general linear groups
over a finite field}, Int. Math. Res. Not. (2000), no.~11,
\href{https://repository.ias.ac.in/36158/1/36158.pdf}{Theorem 2, pp.~581--582}~\cite{prasad}.
At $g=1$ take $H=X-1$, $m_H=d$ and $B_H=1$. The two factors of
$(-1)^{d-1}$ cancel.
\end{proof}

For a squarefree label the fitted character is, by definition, the Hall
product of these primary factors, hence their parabolic induction.
Distinct irreducible polynomials give distinct cuspidal labels (their
degrees differ or their Frobenius orbits differ). Thus the squarefree-label
argument and flag count of \S10.4 apply directly to the fitted character
with the positive sign. This discharges the finite primary-character
identification within F7. The comparison of inner forms and the matching
of archimedean formal degrees in F7 concern different objects.
